\documentclass[10pt]{article}
\usepackage[margin=0.75in]{geometry}
\usepackage{amsmath,amssymb,amsthm,microtype}
\newtheorem{theorem}{Theorem}
\newcommand{\R}{\mathcal R}
\newcommand{\A}{\mathcal A}
\begin{document}
\begin{center}
{\Large An Early Inverse Bridge for Affine-Transfer Induction}\par\medskip
Collatz proof project\quad---\quad September 5, 2026
\end{center}
Write $\R(n)$ for reaching one under the shortcut Collatz map $T$, and
$\A(u)=[\R(3u+2)\Rightarrow\R(27u+20)]$.
The preceding cycle-based construction needed convergence of a base target.
Here an immediate inverse step yields a conditional induction certificate
without that premise. The universal transfer statement remains unproved.

\paragraph{The early predecessor.}
For every natural $w$, put $u=3w+2$ and $v=2w+1$. Then
\[
 v<u,\qquad T(3v+2)=3u+2,\qquad T(27v+20)=81w+71,
\]
whereas $27u+20=81w+74$.
Thus $\R(3u+2)$ and the induction hypothesis $\A(v)$ supply convergence
of $81w+71$. A suitable orbit merge with $81w+74$ completes transfer at $u$.
The source predecessor is available before inspecting any target trajectory.

\begin{theorem}[Uniform early inverse family]
For every natural $Q$, define
\[
 u=59+192Q,\qquad v=39+128Q.
\]
Then $v<u$ and
\begin{align*}
 T(3v+2)&=3u+2,\\
 T^7(27v+20)&=T^6(27u+20)=227+729Q.
\end{align*}
Consequently $\A(v)\Rightarrow\A(u)$.
\end{theorem}
\begin{proof}
The first identity follows by applying the odd shortcut step to
$3v+2=119+384Q$.
For the second, the base paths satisfy
\[
 T^7(1073)=227,\quad j_7(1073)=3,\qquad
 T^6(1613)=227,\quad j_6(1613)=2.
\]
The input increments are $2^7\cdot27Q$ and $2^6\cdot81Q$,
so binary-cylinder transport gives identical endpoint increments $729Q$.
Also $u-v=20+64Q>0$.
Given $\R(3u+2)$, transport convergence backward one step, apply $\A(v)$,
and then use the common endpoint to conclude $\R(27u+20)$.
Lean checks the identities, inequality, and conditional implication for every $Q$.
\end{proof}

\paragraph{Exact search criterion.}
For a binary cylinder $w=r+2^tQ$, let $x,y$ be the $t$-step endpoints of
$81r+71$ and $81r+74$, with odd counts $j,k$.
The endpoint increments are $3^{j+4}Q$ and $3^{k+4}Q$.
Hence $x=y$ and $j=k$ give a uniform merge. The recursive parameters are
\[
 u=3r+2+3\cdot2^tQ,\qquad
 v=2r+1+2^{t+1}Q<u.
\]
These are early certificates: the criterion checks exact finite orbit identities,
not eventual target cycle entry.

\paragraph{Finite evidence.}
At depth 12 the search finds 104 first rules. They cover 407 of the
$4096$ binary residues of $w$, leaving 3689 unresolved. This is coverage
inside $u\equiv2\pmod3$, not of all parameters. Three tests compare complete
small-depth searches with direct orbits, replay all rules on three quotient
lifts, and check the first rule $r=19,t=6$.
The census is Python evidence; the displayed family is kernel checked.

\paragraph{Remaining obligation and verification.}
This family supplies a genuine smaller-parameter induction step without an
unproved target-convergence premise. It does not establish universal rule
coverage. In particular, parameters $2^k-1$ lie in the other two classes
modulo three, so this first inverse template does not resolve the preceding
bootstrap obstruction. Other templates or a global argument are still needed.
\begingroup\raggedright
The three public declarations in \texttt{Collatz.EarlyInverseBridge} enter
the axiom audit. The search is \texttt{analysis/early\_inverse\_bridge\_search.py}.
No mathematical priority claim or proof of Collatz is asserted.
\par\endgroup
\end{document}
